\chapter{非平衡态统计理论初步}

打开隔板后，气体中一小团被标记的分子会向四周扩散。本章从稀薄气体的分子碰撞出发，研究单粒子分布随时间的变化。扩散系数既可以从分子云的展宽求得，也可以从平衡时单个分子的速度记录求得。给分子施加一个弱力，还能测出同一组速度涨落决定的迁移率。

\section{气体分子的碰撞频率}

气体分子通过碰撞使气体达致平衡。
单位时间内碰到单位面积器壁上的粒子数$\Gamma$，定义分布函数$f(\bm r,\bm v,t)$，则
\[
	\Gamma(\bm r,t)=\int_{v_x>0}f(\bm r,\bm v,t)v_x\d\bm v.
\]
这里$x$轴指向器壁，$f$按$\int f\d\bm v=n$归一化。若为平衡态，则满足Maxwell分布$\Gamma=n\ave v/4$。

采用弹性刚球模型（无摩擦，无形变，弹性碰撞）。对于稀薄气体，只考虑两体碰撞，对于两类不同分子间碰撞：单位时间内，平均一个分子1与分子2的碰撞次数$\theta_{12}$，称碰撞频率。

指定速度的分子1的碰撞：
记速度为$\bm v_1$的分子1与第二种分子的平均碰撞频率为$\theta_{12}(\bm r,\bm v_1,t)$，相对速率为$g_{12}=\abs{\bm v_2-\bm v_1}$。两硬球的总截面为$\pi\sigma^2_{12}$，其中$\sigma_{12}=(\sigma_1+\sigma_2)/2$为两分子直径的平均值。给定入射速度后，碰撞参数和碰撞平面决定出射方向。
\begin{align}
	\theta_{12}(\bm r,\bm v_1,t)=\pi\sigma_{12}^2\int f_2(\bm r,\bm v_2,t)g_{12}\d\bm v_2,
\end{align}
以上采用了分子混沌假设，把即将碰撞的两分子的速度联合分布近似为
\[
	f(\bm r,\bm v_1,\bm v_2,t)\simeq f_1(\bm r,\bm v_1,t)f_2(\bm r,\bm v_2,t).
\]

故单位时间内，平均一个分子1与分子2碰撞次数
\begin{align}
	\theta_{12}(\bm r,t)=\frac1{n_1(\bm r,t)}\int f_1(\bm r,\bm v_1,t)\theta_{12}(\bm r,\bm v_1,t)\d\bm v_1.
\end{align}
\paragraph{平均自由程}
设均匀平衡气体由质量$m$、直径$\sigma$相同的硬球组成，数密度为$n$。在气体静止系中随机抽取一对分子，按两份独立的Maxwell分布作平均。速度差的每个分量的方差为$2\kB T/m$，故这对分子的相对速率$g$服从约化质量$m/2$对应的Maxwell分布，$\ave g=\sqrt 2\ave v$，其中$\ave v=\sqrt{8\kB T/(\pi m)}$。若从实际碰撞事件中抽样，速率为$g$的分子对还按碰撞率加权，其相对速率的平均值为$\ave{g^2}/\ave g$。
平均每个分子在单位时间内与其余分子碰撞$\theta=n\pi\sigma^2\ave g$次。用单位时间每个分子走过的平均路程除以平均碰撞次数，得到两次碰撞之间的平均自由程
\begin{equation}
	\ell=\frac{\ave v}{\theta}=\frac{1}{\sqrt2\,n\pi\sigma^2}.
\end{equation}
\paragraph{两体碰撞运动学}
弹性碰撞满足动量守恒和能量守恒
\begin{align}
	m_1\bm v_1+m_2\bm v_2&=m_1\bm v_1'+m_2\bm v_2'\\
	\frac12m_1v_1^2+\frac12m_2v_2^2&=\frac12m_1{v'_1}^2+\frac12m_2{v'_2}^2.
\end{align}
4方程，6未知数，需指定碰撞方向$\bm n$ （或散射角$\theta,\phi$）才能完全决定末态速度。

设$\bm n$为两硬球接触时沿球心连线的单位法向量。无摩擦碰撞的冲量沿此法线，即
\begin{gather*}
	(m_1,\bm v_1)+(m_2,\bm v_2)\quad\overset{\bm n}{\longrightarrow}\quad(m_1,\bm v_1')+(m_2,\bm v_2').\\
	\bm v_1'-\bm v_1=\lambda_1\bm n,\quad\bm v_2'-\bm v_2=\lambda_2\bm n.
\end{gather*}
带入守恒方程
\begin{align*}
	\begin{cases}
		m_1\lambda_1+m_2\lambda_2=0,\\
		m_1\lambda_1(\bm v_1+\bm v_1')\cdot\bm n+m_2\lambda_2(\bm v_2+\bm v_2')\cdot\bm n=0.
	\end{cases}
\end{align*}
解出$\lambda_1,\lambda_2$，
\begin{align}
	\begin{cases}
		\bm v_1'=\bm v_1+\frac{2m_2}{m_1+m_2}\bigb{(\bm v_2-\bm v_1)\cdot\bm n}\bm n\\[1ex]
		\bm v_2'=\bm v_2+\frac{2m_1}{m_1+m_2}\bigb{(\bm v_1-\bm v_2)\cdot\bm n}\bm n
	\end{cases}
\end{align}
对称性：
\begin{compactenum}
	\item 碰撞前后相对速率不变
	\[
		\bm v_1'-\bm v_2'=\bm v_1-\bm v_2-2\bigb{(\bm v_1-\bm v_2)\cdot\bm n}\bm n\implies g_{12}'=g_{12}.
	\]
	\item 碰撞前后相对速度沿碰撞方向变号，垂直方向不变
	\begin{align*}
		(\bm v_1'-\bm v_2')\cdot\bm n&=-(\bm v_1-\bm v_2)\cdot\bm n,\\
		(\bm v_1'-\bm v_2')\times\bm n&=(\bm v_1-\bm v_2)\times\bm n.
	\end{align*}
	\item 原碰撞与逆碰撞等价。所谓逆碰撞：
	\[
		(m_1,\bm v_1')+(m_2,\bm v_2')\quad\overset{-\bm n}{\longrightarrow}\quad(m_1,\bm v_1)+(m_2,\bm v_2).
	\]
\end{compactenum}
\section{Boltzmann输运方程}
讨论分布函数如何随时间变化，记$t$时刻，$(\bm r,\bm v)$处体积元在$\mu$空间$\d\bm r\nd\bm v$中的粒子数$f(\bm r,\bm v,t)\d\bm r\nd\bm v$。

这里使用非简并气体的单粒子分布。对不计内部简并度的气体，相应条件为
\[
	\lambda_T=\Bigp{\frac{h^2}{2\pi m\kB T}}^{1/2}\ll\Bigp{\frac VN}^{1/3}.
\]
硬球的稀薄条件为$n\sigma^3\ll1$，此时平均自由程远大于球径，多体同时碰撞可忽略。下文还假定一次碰撞可用经典散射描述。非简并条件本身只控制全同粒子的交换效应。

因此$\p f/\p t$可看成两部分贡献：漂移项(drift)和碰撞项(collision)
\begin{align}
	\pv ft=\Bigp{\pv ft}_\mathrm{d}+\Bigp{\pv ft}_\mathrm{c}.
\end{align}
漂移项代表运动使$\bm r$变化，外场使$\bm v$变化；碰撞项代表碰撞使$\bm v$变化。
\paragraph{漂移项}首先考虑位置变化：$\d t$内，由$x$处垂直$x$轴的平面进入，和由$x+\d x$处垂直$x$轴的平面离开体积元的粒子数所产生的净增加粒子数
\[
	[f(x,y,\ldots)-f(x+\d x,y,\ldots)]\d y\nd z\nd\bm v\cdot v_x\d t=-v_x\pv fx\d\bm r\nd\bm v\d t.
\]
计所有分量
\[
	-\bm v\cdot\grad_{\bm r}f(\bm r,\bm v,t)\d\bm r\nd\bm v\nd t.
\]

类似的，$x$方向速度变化所产生的净增加粒子数
\[
	-\pp{v_x}(a_xf)\d\bm r\nd\bm v\nd t\implies-\nabla_{\bm v}\cdot(\bm a f)\d\bm r\nd\bm v\nd t.
\]
假设作用力与速度$\bm v$无关，则$\bm a$项可提出来，继而
\begin{align}
	\Bigp{\pv ft}_\mathrm{d}=-\bm v\cdot\grad_{\bm r}f-\bm a\cdot\grad_{\bm v}f.
\end{align}
\paragraph{碰撞项}$\d\bm r\nd\bm v$中分子与别的分子碰撞后离开体积元，称原碰撞，别的分子碰后进入该体积元，称逆碰撞。

原碰撞：$\d t$内，$\d\bm r\nd\bm v$中分子与速度$\bm v_1\to\bm v_1+\d\bm v_1$分子碰撞使粒子数减少
\[
	f(\bm r,\bm v,t)\d\bm r\nd\bm v\cdot f(\bm r,\bm v_1,t)\d\bm v_1\sigma^2\d\Omega\cdot g\d t\cos\theta=f\!f_1\Lambda\d t\nd\Omega\nd\bm r\nd\bm v\nd\bm v_1.
\]
其中$\theta$为入射相对速度与碰撞法线所成的锐角，$\Lambda:=\abs{\bm v-\bm v_1}\sigma^2\cos\theta$。相应的逆碰撞给出
\[
	f'\!f_1'\Lambda\d t\nd\Omega\nd\bm r\nd\bm v\nd\bm v_1.
\]
积分之
\begin{align}
	\Bigp{\pv ft}_\mathrm{c}=\iint(f'\!f_1'-f\!f_1)\Lambda\d\Omega\nd\bm v_1.\quad\theta\in[0,\pi/2].
\end{align}
得到Boltzmann输运方程
\begin{align}
	\pv ft+(\bm v\cdot\grad_{\bm r}+\bm a\cdot\grad_{\bm v})f=\iint(f'\!f_1'-f\!f_1)\Lambda\d\Omega\nd\bm v_1.
\end{align}
方程左边的项即$\d f/\nd t$，这是一个非线性积分微分方程，一般难求解。

精确的单粒子演化方程依赖两粒子分布，后者又依赖三粒子分布，由此形成BBGKY层级\footnote{Bogoliubov--Born--Green--Kirkwood--Yvon hierarchy。}。分子混沌假设忽略碰撞前粒子的速度关联，将入射两粒子分布近似分解为单粒子分布的乘积，从而使Boltzmann方程封闭。这是稀薄气体动力学的近似，不能视为任意时刻的精确独立性。
\section{Boltzmann \textit{H}定理}
为讨论系统趋向平衡时分布函数的变化，定义泛函
\begin{align}
	H(t)=\int f(\bm r,\bm v,t)\ln f(\bm r,\bm v,t)\d\bm r\nd\bm v.
\end{align}
这里$\ln f$表示$\ln(f/f_\ast)$，$f_\ast$是与$f$有相同量纲的常数。粒子数守恒时改变$f_\ast$只给$H$增加一个常数。
\begin{theorem}{Boltzmann $H$定理}{Boltzmann H Theorem}
	若$f(\bm r,\bm v,t)$满足上述Boltzmann输运方程，且空间边界无净熵流（如周期边界或镜面反射边界）、速度无穷远处$f$足够快地衰减，则
	\begin{align}
		\dv Ht\leqslant 0.
	\end{align}
	当且仅当对碰撞积分测度几乎处处有$f\!f_1=f'\!f_1'$时取等号。
\end{theorem}
\begin{proof}
	\begin{align*}
		\dv Ht=\int(1+\ln f)\pv ft\d\bm r\nd\bm v
	\end{align*}
	漂移项贡献为0
	\begin{align*}
		&-\int(1+\ln f)(\bm v\cdot\grad_{\bm r}f+\bm a\cdot\grad_{\bm v}f)\d\bm r\nd\bm v\\
		=&-\int\Bigb{\bm v\cdot\grad_{\bm r}(f\ln f)+\nabla_{\bm v}\cdot(\bm a f\ln f)}\d\bm r\nd\bm v\tag{Gauss}\\
		=&-\int\oint \bm v f\ln f\cdot\d\bm S\nd\bm v-\int\oint\bm af\ln f\cdot\d\bm S_{\bm v}\nd\bm r=0.
	\end{align*}
	碰撞项贡献
	\begin{align*}
		\int (1+\ln f)(f'\!f_1'-f\!f_1)\Lambda\d\Omega\nd\bm v_1\nd\bm r\nd\bm v
	\end{align*}
	交换$\bm v,\bm v_1$，得到的新式与原式相加除2
	\[
		=\frac12\int (2+\ln f\!f_1)(f'\!f_1'-f\!f_1)\Lambda\d\Omega\nd\bm v_1\nd\bm r\nd\bm v
	\]
	又原逆碰撞对称，交换$\bm v,\bm v'$和$\bm v_1,\bm v_1'$，得到
	\[
		=\frac12\int (2+\ln f'\!f_1')(f\!f_1-f'\!f_1')\Lambda\d\Omega\nd\bm v_1\nd\bm r\nd\bm v
	\]
	二者相加除2，得到
	\begin{align}
		\dv Ht=-\frac14\int (\ln f'\!f_1'-\ln f\!f_1)(f'\!f_1'-f\!f_1)\Lambda\d\Omega\nd\bm v_1\nd\bm r\nd\bm v\leqslant 0.
	\end{align}
	被积式中两个差同号，因此取等号要求其乘积对碰撞积分测度几乎处处为零。
\end{proof}

\paragraph{等号条件与局域平衡}
取正的光滑分布$f$。在固定的$(\bm r,t)$，等号条件给出$\ln f+\ln f_1=\ln f'+\ln f_1'$。对弹性硬球的所有碰撞方向都成立的光滑碰撞不变量是$1,\bm v,v^2$的线性组合，分别对应粒子数、动量和能量守恒。因此$\ln f$是速度分量的线性项与$v^2$项之和加常数。配平方并用局域数密度$n$、平均速度$\bm u$和温度$T$表示这些系数，就得到局域Maxwell分布
\[
	f(\bm r,\bm v,t)=n(\bm r,t)\Bigp{\frac{m}{2\pi\kB T(\bm r,t)}}^{3/2}
	\exp\biggb{-\frac{m\abs{\bm v-\bm u(\bm r,t)}^2}{2\kB T(\bm r,t)}}.
\]
这种分布使碰撞项逐点为零。全局平衡还要求分布满足定态漂移方程和边界条件。例如，在无外场的周期盒内，取恒温、$\bm u=0$且初始数密度$n(x)$随位置变化。记归一化的Maxwell速度分布为$M_T(\bm v)$，则初始时刻$\d H/\nd t=0$，但由输运方程$\p f/\p t=-v_x n'(x)M_T(\bm v)$，粒子仍会从不同位置流入和流出。

\paragraph{熵与微观可逆性}
按单粒子概率密度$f/N$取平均，有$H=N\ave{\ln f}$。定义稀薄气体的动理学熵
\begin{equation}
	S=-\kB H+\const.
\end{equation}
在平衡态，这一表达式给出气体的热力学熵。在上述边界条件下，沿Boltzmann方程的解有$\d S/\nd t\geq0$。弛豫所需的时间还取决于碰撞频率和扰动的空间尺度。

对无外场的弹性硬球，把所有分子的速度反向，就能使精确微观运动沿原轨道返回。原先由碰撞产生的出射速度关联，此时成为入射速度关联，因而反演后的气体一般不满足分子混沌假设。这解释了微观运动的可逆性怎样与动理学熵的单调性相容。

Poincaré回归定理也描述微观轨道：当演化保持测度且可达相空间体积有限时，几乎所有初态都会再次进入任意小的初态邻域。对于宏观气体，回归时间通常远长于实验观测时间（参见Huang \S 4.1、\S 4.4、\S 4.5）。

\section{示踪粒子的扩散}
在静止、均匀、温度为$T$的气体中标记少数分子，标记不改变分子质量和碰撞截面。记标记分子的数密度为$c(\bm r,t)$。分子在两次碰撞之间把标记带到邻近位置，反复碰撞使其逐渐失去对初始速度方向的记忆。若$c$在一个平均自由程上的变化很小，标记分布接近$f_0=cM_T(\bm v)$。

用一个与速度无关的弛豫时间$\tau$近似标记分子的碰撞项$-(f-f_0)/\tau$。在远长于$\tau$的扩散时间尺度上，保留浓度梯度的一阶项，输运方程给出
\[
	f-f_0\simeq-\tau\,\bm v\cdot\grad_{\bm r} f_0
	=-\tau M_T(\bm v)\,\bm v\cdot\grad c.
\]
平衡速度分布各向同性，$\int v_i v_jM_T(\bm v)\d\bm v=\delta_{ij}\ave{v^2}/3$。对速度取一阶矩便得到标记分子的数通量和Fick定律
\begin{equation}
	\bm j=\int \bm v f\d\bm v=-D\grad c,\qquad
	D=\frac{\ave{v^2}\tau}{3}=\frac{\kB T}{m}\tau .
\end{equation}
标记分子数在碰撞中守恒。对输运方程按速度积分，得到$\p c/\p t+\div\bm j=0$，因此$c$满足扩散方程$\p c/\p t=D\nabla^2 c$。在三维空间中，若初始标记集中在原点附近且可忽略边界，方程两边乘$r^2$再积分，两次分部积分给出
\[
	\dv{\ave{r^2}}{t}=D\ave{\nabla^2 r^2}=6D.
\]
因此可以测量标记分子云的均方半径随时间增长的斜率来确定$D$。

在量级估算中取$\tau\sim\ell/\ave v$及$\ave{v^2}\sim\ave v^2$，得$D\sim\ave v\ell/3$，量纲为面积除以时间。对给定直径的稀薄硬球，$D$随$\sqrt T/n$增加。一次碰撞保留的速度关联与入射速度、碰撞方向有关，要确定硬球扩散系数的数值前因子，需进一步求解碰撞项。

\section{平衡涨落与线性响应}\label{sec:linear-response}

\subsection{从速度记录求扩散系数}
在均匀、各向同性的三维经典流体中跟踪一个粒子，取流体静止系，使$\ave{\bm v}=0$。速度自关联函数
\[
	C_v(t)=\ave{\bm v(t)\cdot\bm v(0)}
\]
度量粒子保留初始运动方向的程度。平均可以取遍许多平衡轨迹，也可以在遍历的平衡轨迹上取不同的时间原点。每条轨迹都满足$\Delta\bm r(t)=\int_0^t\bm v(s)\d s$。把位移平方，利用平衡态的时间平移不变性以及$C_v(-s)=C_v(s)$，得到
\begin{equation}
	\ave{\abs{\Delta\bm r(t)}^2}
	=\int_0^t\d s\int_0^t\d s'\,C_v(s-s')
	=2\int_0^t(t-s)C_v(s)\d s.
	\label{eq:sm-msd-vacf}
\end{equation}
短时间内$C_v(s)\simeq\ave{v^2}$，均方位移按$t^2$增长，此时可以分辨粒子沿初始方向的飞行。若速度关联可积，长时间均方位移的斜率趋向$6D$。\eqref{eq:sm-msd-vacf}的导数给出自扩散的Green--Kubo公式
\begin{equation}
	D=\frac13\int_0^\infty C_v(s)\d s
	=\int_0^\infty\ave{v_x(s)v_x(0)}\d s .
	\label{eq:sm-green-kubo-d}
\end{equation}
在分子动力学模拟中，均方位移的长时斜率和速度关联积分的平台都可以用来确定$D$。若速度关联按时间常数$\tau$指数衰减，其积分为$C_v(0)\tau$，代入便给出上节的扩散系数。真实流体中的粒子也可能在碰到邻居后折返，使某段时间内的速度关联为负，这一段会减小扩散系数。

\subsection{弱力、迁移率与Einstein关系}
沿$x$方向给示踪粒子一个很小的冲量$I$。设该粒子的动量只通过$p_x^2/(2m)$进入经典哈密顿量，冲量把$p_x$变为$p_x+I$。脉冲后的相空间分布用脉冲前的正则分布$\rho_0$表示为
\[
	\rho_{0^+}(\ldots,p_x)=\rho_0(\ldots,p_x-I)
	=\rho_0(\ldots,p_x)\bigb{1+\beta I v_x+O(I^2)},\qquad
	\beta=\frac1{\kB T}.
\]
此后粒子仍按原哈密顿量运动。在上述分布中对$v_x(t)$取平均，得到脉冲响应$\delta\ave{v_x(t)}=\beta I\ave{v_x(t)v_x(0)}$至$I$的一阶。把任意弱力$F(t')$分解为连续的小冲量$F(t')\d t'$，便有
\[
	\delta\ave{v_x(t)}
	=\beta\int_{-\infty}^{t}\ave{v_x(t-t')v_x(0)}F(t')\d t'.
\]
若恒力作用足够久，粒子取得稳定的漂移速度，定义迁移率$b$为漂移速度与弱力之比。与\eqref{eq:sm-green-kubo-d}比较得
\begin{equation}
	b=\beta\int_0^\infty\ave{v_x(s)v_x(0)}\d s,\qquad
	D=\kB T b .
	\label{eq:sm-einstein-mobility}
\end{equation}
这就是Einstein关系。举例来说，若长时间的阻力为$-\zeta v_x$，力平衡给出$b=1/\zeta$，于是平衡粒子的扩散系数为$\kB T/\zeta$。将速度关联除以初值$\ave{v_x^2}=\kB T/m$再积分，得到速度记忆的时间尺度$m/\zeta$。

\subsection{量子Kubo公式与涨落耗散定理}
设量子系统起初处于$\hat H_0$的正则平衡态$\hat\rho_0$，外场$f(t)$通过$\hat H'(t)=-f(t)\hat B$耦合到可观测量$\hat B$。系统对这个场的响应可以通过另一个可观测量$\hat A$来测定。在相互作用绘景中，将密度算符的运动方程积分至$f$的一阶得
\[
	\delta\hat\rho_I(t)=\frac{\i}{\hbar}\int_{-\infty}^t
	f(t')\bigb{\hat B(t'),\hat\rho_0}\d t'.
\]
这里算符均按$\hat H_0$演化。把上式乘以$\hat A(t)$取迹，利用迹的循环性得到Kubo公式\footnote{参见David Tong的\href{https://davidtong.org/pdfs/teaching/kinetic-theory/kinetic4.pdf}{\textit{Kinetic Theory}第四章}。}
\begin{align}
	\delta\ave{\hat A(t)}&=\int_{-\infty}^{\infty}
	\chi^R_{AB}(t-t')f(t')\d t',\notag\\
	\chi^R_{AB}(s)&=\frac{\i}{\hbar}\Theta(s)
	\ave{[\hat A(s),\hat B(0)]}_0 .
	\label{eq:sm-kubo-retarded}
\end{align}
阶跃函数$\Theta(s)$体现因果性，外场只对其施加时刻之后的测量有贡献。

取傅里叶变换$g(\omega)=\int_{-\infty}^{\infty}\e^{\i\omega t}g(t)\d t$。设$\hat A=\hat B$为厄米算符，记$\chi''_{AA}(\omega)=\Im\chi^R_{AA}(\omega)$。施加$f(t)=f_0\cos\omega t$时，外场的周期平均输入功率为
\[
	\avg{P}=-\avg{\dot f\,\ave{\hat A}}
	=\tfrac12\omega f_0^2\chi''_{AA}(\omega).
\]
从而，测量吸收功率即可得到响应的耗散部分。平衡中也能测到同一组跃迁。令$\delta\hat A=\hat A-\ave{\hat A}_0$，定义未对称化的涨落谱
\[
	S_A(\omega)=\int_{-\infty}^{\infty}\e^{\i\omega t}
	\ave{\delta\hat A(t)\delta\hat A(0)}_0\d t .
\]
在$\hat H_0$本征态中插入恒等算符，记$p_m=\e^{-\beta E_m}/Z$、$\omega_{nm}=(E_n-E_m)/\hbar$，即得
\[
	S_A(\omega)=2\pi\sum_{m,n}p_m\abs{\mel{m}{\delta\hat A}{n}}^2
	\delta(\omega-\omega_{nm}).
\]
交换$m,n$并用$p_n/p_m=\e^{-\beta\hbar\omega_{nm}}$，可见正负频率的权重满足细致平衡$S_A(-\omega)=\e^{-\beta\hbar\omega}S_A(\omega)$。另一方面，\eqref{eq:sm-kubo-retarded}中的对易子给出$2\hbar\chi''_{AA}(\omega)=S_A(\omega)-S_A(-\omega)$。消去两种跃迁权重之差，得到非零频率处的量子涨落耗散定理
\begin{equation}
	S_A^{\mathrm{sym}}(\omega)
	:=\frac{S_A(\omega)+S_A(-\omega)}2
	=\hbar\coth\!\Bigp{\frac{\beta\hbar\omega}{2}}\chi''_{AA}(\omega).
	\label{eq:sm-quantum-fdt}
\end{equation}
左边是对称化的平衡噪声谱。若$\hbar\abs{\omega}\ll\kB T$，便得到经典关系$S_A^{\mathrm{sym}}(\omega)\simeq 2\kB T\chi''_{AA}(\omega)/\omega$。在相反的极限$\hbar\omega\gg\kB T$下，$S_A(-\omega)$的权重受到指数抑制，因为能向外场释放能量$\hbar\omega$的激发态很少被占据。
